%============================================================
% CH1 · PATTERN 4 — IMPLICIT STATE GRAPHS: THEORY
%============================================================
\subsection{Pattern Theory}
\label{ssec:ch1-p4-theory}

\begin{patternbox}{Pattern 4 at a Glance}
The input does \textbf{not} contain a graph. It describes a \textbf{situation} (a number on a display, the digits of a lock, the water in two jugs) and the \textbf{moves} that change it. Every possible situation is a \textbf{vertex} (a \emph{state}), every move is an \textbf{edge}. ``Fewest moves'' $=$ shortest path, and because every move costs one, \textbf{BFS} finds it --- generating the neighbours of a state on the fly instead of storing the graph.

\begin{center}
\begin{tikzpicture}[x=2.3cm, y=0.95cm]
    \node[font=\scriptsize\bfseries, text=darkgray] at (0,1.15) {dist 0};
    \node[font=\scriptsize\bfseries, text=darkgray] at (1,1.15) {dist 1};
    \node[font=\scriptsize\bfseries, text=darkgray] at (2,1.15) {dist 2};
    \node[font=\scriptsize\bfseries, text=darkgray] at (3,1.15) {dist 3};
    \node[gvsrc] (s3) at (0,-1.5) {3};
    \node[gvdone] (s6) at (1,0) {6};
    \node[gvdone] (s2) at (1,-3) {2};
    \node[gvdone] (s12) at (2,0.2) {12};
    \node[gvdone] (s5) at (2,-1.0) {5};
    \node[gvdone] (s4) at (2,-2.2) {4};
    \node[gvdone] (s1) at (2,-3.4) {1};
    \node[gvseen] (s11) at (3,0.2) {11};
    \node[gvdst] (s10) at (3,-1.0) {10};
    \node[gvseen] (s8) at (3,-2.2) {8};
    \draw[gtree, -{Stealth[length=2.4mm]}] (s3) -- node[gw, above, sloped]{$\times2$} (s6);
    \draw[gtree, -{Stealth[length=2.4mm]}] (s3) -- node[gw, below, sloped]{$-1$} (s2);
    \draw[gtree, -{Stealth[length=2.4mm]}] (s6) -- (s12);
    \draw[gtree, -{Stealth[length=2.4mm]}] (s6) -- (s5);
    \draw[gtree, -{Stealth[length=2.4mm]}] (s2) -- (s4);
    \draw[gtree, -{Stealth[length=2.4mm]}] (s2) -- (s1);
    \draw[gtree, -{Stealth[length=2.4mm]}] (s12) -- (s11);
    \draw[gpath, -{Stealth[length=2.6mm]}] (s5) -- node[gw, above]{$\times2$} (s10);
    \draw[gtree, -{Stealth[length=2.4mm]}] (s4) -- (s8);
    \draw[gfaded, dashed, -{Stealth[length=2mm]}] (s5) -- (s4);
    \node[note, text width=3.3cm, anchor=west] at (3.35,-1.9) {moves: $\times2$ and $-1$\\start $3$, goal $10$\\answer $=$ layer of the goal $=3$\\(dashed: $5\to4$ finds an\\already discovered state)};
\end{tikzpicture}
\end{center}

\begin{lstlisting}[style=snippet]
dist[start] = 0;  queue = {start};               // layer 0
while (!queue.empty()) {
    s = pop_front();  if (s == goal) return dist[s];
    for (each move m)  { t = m(s);                // neighbour generated, not stored
        if (valid(t) && dist[t] == -1) { dist[t] = dist[s] + 1; push_back(t); } }
}
return -1;                                       // goal not reachable
\end{lstlisting}
Time $\bigO{\#\text{states} \times \#\text{moves}}$ --- so the first job is always to \textbf{bound the number of states}.
\end{patternbox}

%------------------------------------------------------------
\subsubsection*{From Pattern 3 to Pattern 4}

In Pattern~3 the graph was hidden in a grid, but at least you could \emph{see} it: the cells were in the input and the neighbours were the four cells around. Here even that disappears. In Two Buttons (Problem~1.4.1) the input is just two numbers, $n$ and $m$; in Open the Lock (Problem~1.4.2) it is a target code and a list of forbidden codes. Nobody hands you vertices or edges --- you have to \emph{invent} them.

\begin{diagrambox}
\begin{tikzpicture}
    % left: explicit graph
    \node[font=\small\bfseries, text=deepblue] at (1.5,2.2) {Patterns 1--3: the graph is given};
    \node[gv] (a) at (0,1) {1};
    \node[gv] (b) at (1.5,1.5) {2};
    \node[gv] (c) at (3,1) {3};
    \node[gv] (d) at (1.5,0) {4};
    \draw[ge] (a)--(b) (b)--(c) (a)--(d) (d)--(c);
    \node[draw=darkgray!60, rounded corners=2pt, font=\scriptsize\ttfamily, align=left, fill=lightgray, anchor=north]
        at (1.5,-0.6) {1 2\\2 3\\1 4\\4 3};
    \node[note] at (1.5,-2.3) {edges are read from the input\\and stored in \texttt{adj}};
    % right: implicit graph
    \node[font=\small\bfseries, text=deepblue] at (8.2,2.2) {Pattern 4: the graph is a \emph{rule}};
    \node[gvcur, minimum size=9mm] (s) at (6.4,0.5) {$s$};
    \node[tag=accentorange, text width=2.5cm] (f) at (8.4,0.5) {\scriptsize\texttt{neighbours(s)}\\[-1pt]\scriptsize apply every move};
    \node[gvseen, minimum size=9mm, font=\scriptsize\bfseries] (t1) at (10.6,1.5) {$2s$};
    \node[gvseen, minimum size=9mm, font=\scriptsize\bfseries] (t2) at (10.6,-0.5) {$s{-}1$};
    \draw[flow] (s) -- (f);
    \draw[flow] (f) -- (t1);
    \draw[flow] (f) -- (t2);
    \node[note] at (8.4,-1.6) {nothing is stored: the edges of $s$\\are \emph{computed} when BFS reaches $s$};
\end{tikzpicture}
\figcap{An explicit graph is data; an implicit graph is a function. BFS only ever asks ``what are the neighbours of this vertex?'' --- and a function can answer that as well as an adjacency list.}
\end{diagrambox}

This is the most important modelling skill of the whole book: whenever a problem says ``minimum number of operations / moves / presses / steps to turn $A$ into $B$'', ask yourself \emph{what the situation between two moves looks like}. That snapshot is a vertex.

%------------------------------------------------------------
\subsubsection*{The three questions}

Before writing any code, answer three questions on paper. They are the Graph Model of Pattern~4.

\begin{diagrambox}
\begin{tikzpicture}[node distance=5mm]
    \node[fcproc=3.3cm] (q1) {\textbf{1. State}\\what must I know\\between two moves?};
    \node[fcproc=3.3cm, right=of q1] (q2) {\textbf{2. Moves}\\which states can one\\move reach? (edges)};
    \node[fcproc=3.3cm, right=of q2] (q3) {\textbf{3. Size}\\how many states\\can ever matter?};
    \node[fcout=2.2cm, right=of q3] (bfs) {BFS\\from start};
    \draw[fcarrow] (q1) -- (q2);
    \draw[fcarrow] (q2) -- (q3);
    \draw[fcarrow] (q3) -- (bfs);
    \node[note, below=2mm of q1, text width=3.3cm] {display: one number\\lock: 4 digits\\jugs: a pair $(a,b)$};
    \node[note, below=2mm of q2, text width=3.3cm] {$\times2$, $-1$\\turn one wheel $\pm1$\\fill / empty / pour};
    \node[note, below=2mm of q3, text width=3.3cm] {$\le 2\cdot10^4$\\$10^4$\\$(x{+}1)(y{+}1)$};
\end{tikzpicture}
\figcap{The three questions, with the answers for the two solved problems and for the water jugs of Exercise~E4.3.}
\end{diagrambox}

\begin{enumerate}[leftmargin=*]
    \item \textbf{State} --- the \emph{smallest} description of the situation that decides which moves are possible and whether we have arrived. Too little information and two different situations get merged (wrong answers); too much and the state space explodes (TLE). The \emph{history} of how you got there is \textbf{not} part of the state --- BFS remembers the number of moves for you.
    \item \textbf{Moves} --- every move of the statement becomes one line of code that turns $s$ into a neighbour $t$, plus a validity test (``the display must stay positive'', ``a dead-end code may not be entered'').
    \item \textbf{Size} --- multiply the ranges of the components (4 digits: $10^4$; two jugs: $(x+1)(y+1)$). If a component is unbounded (a number that can grow forever), you must \emph{prove} a window outside which states are useless. Then $\#\text{states} \times \#\text{moves}$ must fit in the time limit (about $10^8$ simple steps per second).
\end{enumerate}

%------------------------------------------------------------
\subsubsection*{Turning a state into an array index}

BFS needs a \texttt{dist} entry for each state. If the state is a small integer, it \emph{is} the index. Otherwise \textbf{encode} it as an integer --- a mixed-radix number --- so that \texttt{dist} can be a plain vector instead of a slow \texttt{map}.

\begin{diagrambox}
\begin{tikzpicture}
    % lock code
    \node[font=\small\bfseries, text=deepblue] at (1.4,1.15) {lock wheels};
    \foreach \d [count=\i from 0] in {0,2,0,2}
        \node[gcell, minimum size=8mm, font=\large\bfseries] (w\i) at (\i*0.85,0) {\d};
    \foreach \p [count=\i from 0] in {1000,100,10,1}
        \node[font=\scriptsize, text=darkgray] at (\i*0.85,-0.65) {$\times\p$};
    \draw[flow] (3.2,0) -- (4.2,0);
    \node[tag=successgreen] at (5.35,0) {\texttt{id} $=202$};
    \node[note] at (2.6,-1.35) {4 digits $\Rightarrow$ ids $0..9999$};
    % jug pair
    \node[font=\small\bfseries, text=deepblue] at (9.3,1.15) {pair $(a,b)$, $0\le b\le y$};
    \node[tag=primaryblue] (p) at (8.3,0) {$(a,b)=(2,4)$};
    \draw[flow] (9.4,0) -- (10.3,0);
    \node[tag=successgreen, align=center] at (11.75,0) {\texttt{id} $= a\cdot(y{+}1)+b$\\$=2\cdot6+4=16$};
    \node[note] at (10.2,-1.35) {$y=5$: ids $0..(x{+}1)(y{+}1){-}1$; decode:\\$a=\texttt{id}/(y{+}1)$, \ $b=\texttt{id}\bmod(y{+}1)$};
\end{tikzpicture}
\figcap{Two encodings used in this pattern. Any tuple $(c_1,\dots,c_k)$ with $0 \le c_i < R_i$ becomes $(\dots(c_1 R_2 + c_2)R_3 + \dots)R_k + c_k$; a tuple of booleans becomes a bitmask.}
\end{diagrambox}

%------------------------------------------------------------
\subsubsection*{Why BFS gives the fewest moves}

BFS explores the states in \textbf{layers}: first the start (0 moves), then everything one move away, then everything two moves away, and so on --- like ripples on water. The first time BFS discovers a state is therefore along a shortest sequence of moves.

\begin{diagrambox}
\begin{tikzpicture}
    \foreach \r/\c in {3.0/primaryblue!6, 2.1/primaryblue!12, 1.2/primaryblue!20}
        \fill[\c] (0,0) circle (\r);
    \foreach \r in {1.2, 2.1, 3.0} \draw[primaryblue!50, dashed] (0,0) circle (\r);
    \node[gvsrc] at (0,0) {$s$};
    \node[font=\scriptsize\bfseries, text=primaryblue] at (0,0.85) {dist 1};
    \node[font=\scriptsize\bfseries, text=primaryblue] at (0,1.75) {dist 2};
    \node[font=\scriptsize\bfseries, text=primaryblue] at (0,2.65) {dist 3};
    \node[gvdst, minimum size=6mm] (g) at (1.9,-1.6) {$g$};
    % queue picture
    \begin{scope}[xshift=4.6cm, yshift=0.3cm]
        \node[font=\small\bfseries, text=deepblue, anchor=west] at (0,1.2) {the queue at any moment};
        \foreach \v/\s [count=\i from 0] in {d/gcell, d/gcell, d/gcell, d{+}1/gcell, d{+}1/gcell}
            \node[\s, minimum width=1.1cm, font=\small] (q\i) at (\i*1.1+0.55,0.3) {$\v$};
        \node[font=\scriptsize, text=darkgray] at (0.55,-0.4) {front};
        \node[font=\scriptsize, text=darkgray] at (4.95,-0.4) {back};
        \node[note, anchor=west, text width=6cm, align=left] at (0,-1.5)
            {Distances in the queue are \emph{sorted} and differ by at most 1: some states at distance $d$, followed by some at $d+1$. Popping from the front always takes a closest unfinished state.};
    \end{scope}
\end{tikzpicture}
\figcap{Left: layers of a BFS. Right: the queue invariant that makes the layers come out in order.}
\end{diagrambox}

\begin{invariantbox}[{Invariant of BFS with unit-cost moves}]
\textbf{(I1)} At every moment the queue holds states with distances $d, \dots, d, d{+}1, \dots, d{+}1$ (non-decreasing, spread at most 1).
\textbf{(I2)} When a state $t$ gets $\texttt{dist}[t] = \texttt{dist}[s]+1$, this is the true minimum number of moves to $t$.

\emph{Proof sketch.} (I1) holds for the queue $\{\text{start}\}$. Popping a state at distance $d$ and pushing its new neighbours at $d{+}1$ to the back keeps the order and the spread. For (I2): suppose $t$ could be reached in $k < \texttt{dist}[s]+1$ moves, and let $p$ be the state just before $t$ on such a shorter route; $p$ is at distance $k-1 < \texttt{dist}[s]$. By (I1) $p$ was popped \emph{before} $s$, and at that moment it would have discovered $t$ first, with the smaller value --- contradiction. \hfill$\blacksquare$

Consequence: mark a state as discovered \textbf{when it is pushed}, not when it is popped; and the first time the goal is popped (or pushed), its distance is the answer.
\end{invariantbox}

%------------------------------------------------------------
\subsubsection*{Bounding an infinite state space}

A number that can be doubled forever gives infinitely many states, and BFS would never stop when the goal is unreachable. The fix is never ``pick a big constant and hope'': it is a short \textbf{proof} that leaving some window $[\mathrm{LO}, \mathrm{HI}]$ cannot help.

\begin{diagrambox}
\begin{tikzpicture}[x=1cm]
    \fill[recursionbg] (-0.4,-0.25) rectangle (1,0.25);
    \fill[solutionbg] (1,-0.25) rectangle (9,0.25);
    \fill[recursionbg] (9,-0.25) rectangle (12.4,0.25);
    \draw[thick, -{Stealth}] (-0.4,0) -- (12.5,0);
    \foreach \x/\l in {1/LO, 9/HI} {
        \draw[thick] (\x,-0.35) -- (\x,0.35);
        \node[font=\small\bfseries, text=deepblue] at (\x,0.65) {\l};
    }
    \node[font=\scriptsize, text=warningred, align=center] at (0.3,-0.75) {invalid\\(e.g.\ $\le 0$)};
    \node[font=\scriptsize, text=successgreen!70!black, align=center] at (5,-0.75) {every state an optimal solution can visit};
    \node[font=\scriptsize, text=warningred, align=center] at (10.7,-0.75) {provably useless\\(``too far to come back cheaply'')};
    \node[gvsrc, minimum size=6mm] at (3,0) {};
    \node[gvdst, minimum size=6mm] at (6.5,0) {};
\end{tikzpicture}
\figcap{BFS only on the green window. Two Buttons: HI $=\max(n, 2m)$ (Lemma in Problem~1.4.1). Exercise E4.5: only values in $[0,1000]$ can be continued, everything else is a leaf.}
\end{diagrambox}

Typical arguments:
\begin{itemize}[leftmargin=*]
    \item \textbf{Only-increasing moves} ($+3$, $\times2$, appending a digit): a value above the goal can never come back, so HI $=$ goal.
    \item \textbf{Overshoot costs more than it saves}: in Two Buttons, a doubling made at a value $\ge m$ adds at least $m$ that must later be removed one by one, so it is never optimal; hence HI $= \max(n, 2m)$.
    \item \textbf{Modular states}: ``$\bmod\ 32768$'' or a wheel with 10 digits makes the space finite automatically.
    \item \textbf{The statement does it for you}: ``operations are possible only while $0\le x\le 1000$''.
\end{itemize}

%------------------------------------------------------------
\subsubsection*{When to recognise it}

\begin{itemize}[leftmargin=*]
    \item ``\textbf{minimum number} of operations / moves / presses / turns / steps'' to go from a start configuration to a goal;
    \item a \textbf{small menu of moves}, each costing the same (one press, one turn, one pour);
    \item the configuration fits in a few small numbers, so the number of states is at most a few million;
    \item ``is it \textbf{possible} at all?'' --- the same BFS (or a DFS) answers reachability.
\end{itemize}
\textbf{Not this pattern (or not only):}
\begin{itemize}[leftmargin=*]
    \item moves have \textbf{different costs} --- BFS is wrong; use 0-1 BFS or Dijkstra (Chapters 7--8);
    \item the number of states is astronomically large (values up to $10^{18}$, permutations of 20 items) --- look for a mathematical shortcut or a greedy in the \emph{reverse} direction (Problem~1.4.1, Approach~3);
    \item you must \textbf{count} the ways or optimise a sum over a path in a graph without cycles --- that is DP on a DAG (Chapter 5 and the DP book);
    \item a single long walk where each state has exactly one successor --- that is a functional graph (Pattern~5).
\end{itemize}

%------------------------------------------------------------
\subsubsection*{Variations map}

\begin{diagrambox}
\begin{tikzpicture}[
    var/.style={rectangle, rounded corners=2pt, draw=patternborder, fill=patternbg,
                font=\scriptsize, align=center, text width=3.25cm, minimum height=1.25cm, inner sep=3pt}]
    \node[var, fill=answerfill, draw=accentorange, text width=4.6cm, font=\scriptsize\bfseries] (root) at (0,0)
        {BFS over states\\fewest unit-cost moves from start to goal};
    \node[var] (v1) at (-5.4,-2) {\textbf{numbers}\\(Two Buttons, E4.1, E4.4, E4.5)\\bound the window first};
    \node[var] (v2) at (-1.8,-2) {\textbf{digit tuples / codes}\\(Open the Lock)\\mixed-radix id};
    \node[var] (v3) at (1.8,-2)  {\textbf{resource pairs}\\(water jugs, E4.3)\\$(a,b)$ as $a(y{+}1)+b$};
    \node[var] (v4) at (5.4,-2)  {\textbf{strings / boards}\\(word ladders, sliding puzzles)\\Chapter 3};
    \node[var] (v5) at (-3.6,-3.6) {\textbf{reverse the moves}\\multi-source / one target\\(E4.4, greedy of 1.4.1)};
    \node[var] (v6) at (0,-3.6)    {\textbf{meet in the middle}\\bidirectional BFS\\(Problem 1.4.2, Approach 3)};
    \node[var] (v7) at (3.6,-3.6)  {\textbf{weighted moves}\\0-1 BFS / Dijkstra\\(Chapters 7--8)};
    \node[draw=patternborder, dashed, rounded corners=4pt, fit=(v1)(v4)(v5)(v7), inner sep=5pt] (grp) {};
    \draw[flow] (root.south) -- (grp.north);
\end{tikzpicture}
\figcap{Top row: what a state looks like. Bottom row: how the search itself changes.}
\end{diagrambox}

\textbf{Twists and how to handle them}
\begin{itemize}[leftmargin=*]
    \item \textbf{Forbidden states} (dead ends, walls, ``the device breaks''): simply never push them. If the \emph{start} is forbidden, answer immediately.
    \item \textbf{Many queries, one goal}: run one BFS \emph{backwards} from the goal over reversed moves; then every start is answered by a lookup (Exercise~E4.4).
    \item \textbf{Print the sequence of moves}: store \texttt{par[t] = s} when $t$ is discovered and walk back from the goal (Template~1).
    \item \textbf{Goal described by a property} (``total water equals $T$''): test the property when a state is popped instead of comparing with one goal state.
    \item \textbf{Out-of-window results that are still allowed} (Exercise~E4.5): check them against the goal, but do not push them.
\end{itemize}

%------------------------------------------------------------
\subsubsection*{Code templates}

\textbf{Template 1 --- integer states in a window, with the path.} Moves of the demo: $+3$ and $\times2$. Both only increase the value, so the window is $[\text{start}, \text{goal}]$.

\begin{lstlisting}[style=cpp]
#include <bits/stdc++.h>
using namespace std;

// fewest moves from start to goal inside the window [LO, HI]; path receives the states
int bfsInt(int start, int goal, int LO, int HI, vector<int>& path) {
    vector<int> dist(HI - LO + 1, -1), par(HI - LO + 1, -1);
    queue<int> q;
    dist[start - LO] = 0;
    q.push(start);
    while (!q.empty()) {
        int s = q.front(); q.pop();
        if (s == goal) break;                          // goal popped: distance final
        int next[2] = {s + 3, s * 2};                  // the moves of this demo
        for (int t : next) {
            if (t < LO || t > HI) continue;            // outside the proven window
            if (dist[t - LO] != -1) continue;          // already discovered
            dist[t - LO] = dist[s - LO] + 1;           // mark when PUSHED
            par[t - LO] = s;
            q.push(t);
        }
    }
    path.clear();
    if (dist[goal - LO] == -1) return -1;
    for (int v = goal; v != -1; v = par[v - LO]) path.push_back(v);
    reverse(path.begin(), path.end());
    return dist[goal - LO];
}

int main() {
    vector<int> path;
    cout << bfsInt(1, 25, 1, 25, path) << ":";
    for (int v : path) cout << " " << v;               // 5: 1 4 8 11 22 25
    cout << "\n";
}
\end{lstlisting}

\textbf{Template 2 --- a tuple of booleans as a bitmask.} Five lamps in a row; pressing button $i$ toggles lamp $i$ and its neighbours. Fewest presses from ``all off'' to a goal pattern. There are $2^5 = 32$ states and 5 moves each.

\begin{lstlisting}[style=cpp]
#include <bits/stdc++.h>
using namespace std;

int main() {
    const int N = 5, S = 1 << N;                       // state = bitmask of lamps that are on
    vector<int> dist(S, -1);
    queue<int> q;
    dist[0] = 0;
    q.push(0);
    while (!q.empty()) {
        int s = q.front(); q.pop();
        for (int i = 0; i < N; i++) {
            int mask = 1 << i;                         // lamp i ...
            if (i > 0) mask |= 1 << (i - 1);           // ... and its neighbours
            if (i + 1 < N) mask |= 1 << (i + 1);
            int t = s ^ mask;
            if (dist[t] == -1) { dist[t] = dist[s] + 1; q.push(t); }
        }
    }
    int reachable = 0;
    for (int d : dist) if (d != -1) reachable++;
    cout << dist[S - 1] << " " << dist[1] << " " << reachable << "\n";   // 2 -1 16
}
\end{lstlisting}
All five lamps on needs 2 presses (buttons 0 and 3); lamp 0 alone (\texttt{dist[1]}) is \emph{impossible}: only 16 of the 32 patterns are reachable. Reachability questions are answered by the same BFS.

\textbf{Template 3 --- string states with a hash map.} Fewest swaps of adjacent letters turning one word into another. Convenient when the state has no obvious small encoding; slower by a constant factor, so prefer an integer id when the state space is large.

\begin{lstlisting}[style=cpp]
#include <bits/stdc++.h>
using namespace std;

int fewestSwaps(const string& from, const string& to) {
    unordered_map<string, int> dist;                   // state -> number of moves
    queue<string> q;
    dist[from] = 0;
    q.push(from);
    while (!q.empty()) {
        string s = q.front(); q.pop();
        if (s == to) return dist[s];
        for (size_t i = 0; i + 1 < s.size(); i++) {
            string t = s;
            swap(t[i], t[i + 1]);                      // one move = one adjacent swap
            if (!dist.count(t)) { dist[t] = dist[s] + 1; q.push(t); }
        }
    }
    return -1;
}

int main() {
    cout << fewestSwaps("aabb", "bbaa") << " " << fewestSwaps("abcd", "dcba") << "\n";  // 4 6
}
\end{lstlisting}

%------------------------------------------------------------
\subsubsection*{Common mistakes}

\begin{warningbox}[Mistake 1 --- no bound on the state space]
\begin{lstlisting}[style=snippet]
q.push(2 * s);                               // WRONG: values grow forever, unreachable goal = infinite loop
if (2 * s <= HI) q.push(2 * s);              // RIGHT: HI comes from a proof, not from a guess
\end{lstlisting}
A window that is too \emph{small} is just as bad: in Two Buttons with HI $= m$ the answer for $n=3$, $m=5$ (path $3\to6\to5$) is not found at all.
\end{warningbox}

\begin{warningbox}[Mistake 2 --- marking visited when popping]
If \texttt{dist[t]} is set only when $t$ is popped, the same state is pushed many times (once per incoming edge). The answer stays correct, but the queue can grow to $\#\text{states}\times\#\text{moves}$ entries --- \mle{} or \tle{} on large spaces. Mark on push.
\end{warningbox}

\begin{warningbox}[Mistake 3 --- DFS for ``fewest moves'']
A DFS with a visited array finds \emph{a} way, usually a long one: on the lock it may wander through thousands of codes. Only BFS (unit costs) guarantees the minimum.
\end{warningbox}

\begin{warningbox}[Mistake 4 --- a state that forgets something]
If the moves or the goal test depend on a quantity, it must be in the state. Example: with a limited number of ``special'' moves, the state is (position, special moves left), not just the position --- otherwise BFS merges situations that behave differently.
\end{warningbox}

\begin{warningbox}[Mistake 5 --- forgetting the start or the goal edge cases]
Start $=$ goal $\Rightarrow$ 0 moves. Start forbidden $\Rightarrow$ $-1$ before the loop (Open the Lock with \texttt{"0000"} among the dead ends). Goal outside the window $\Rightarrow$ check it explicitly.
\end{warningbox}

\begin{checklistbox}[``I'm stuck'' checklist for Pattern 4]
\begin{itemize}[leftmargin=*]
    \item[$\square$] What is the smallest snapshot that decides the next moves and the goal test? That is the state.
    \item[$\square$] Did I write down every move of the statement, including wrap-around and ``until full / until empty''?
    \item[$\square$] How many states are there? Is there a \emph{proof} for my window? Does $\#\text{states}\times\#\text{moves}$ fit?
    \item[$\square$] Can I encode the state as an integer for a flat \texttt{dist} array?
    \item[$\square$] Do all moves cost the same? (Otherwise: not plain BFS.)
    \item[$\square$] Many starts and one goal --- should I run BFS backwards from the goal?
    \item[$\square$] Is the state space too big? Try the reverse direction, a greedy, or meet in the middle.
\end{itemize}
\end{checklistbox}
